\subsubsection{Shao and Wu: Moment Contraction for Nonlinear Time Series}
\label{sec:shao-wu-gmc}

\paragraph{Published result.}
\citet{shao2007spectral} develop spectral theory for nonlinear time series.
Section~5, Theorem~5.1 (pp.~11--12 of the linked electronic reprint), provides
the relevant generation result. With i.i.d. innovations and a measurable,
time-homogeneous recurrence $\mathcal R_{\mathcal T}(\mathbf{s};\varepsilon)$,
its second-moment specialization assumes a finite second moment at some
reference history and the global bound
\begin{equation}
\begin{split}
 \|\mathcal R_{\mathcal T}(\mathbf{s};\varepsilon)
       -\mathcal R_{\mathcal T}(\mathbf{s}';\varepsilon)\|_{L^2}
 &\leq \sum_{d=0}^{L-1}a_d^{\mathrm{GMC}}|s_d-s'_d|,\\
 a_d^{\mathrm{GMC}}\geq0,\qquad
 \sum_{d=0}^{L-1}a_d^{\mathrm{GMC}}&<1.
\end{split}
\label{eq:gmc-generator}
\end{equation}
Here $\|Z\|_{L^2}=(\mathbb E[\|Z\|_2^2])^{1/2}$. The theorem yields a
causal strictly stationary solution and geometric moment contraction (GMC):
the expected squared difference between coupled versions with a shared future
innovation sequence decays geometrically. Vector-valued states are admitted;
the absolute differences become norms of lag blocks. Theorem~5.2 preserves
GMC under stable ARMA filtering, and Examples~5.1--5.5 cover exponential
autoregression, conditional heteroscedasticity, and bilinear recursions.

\paragraph{Proposed generator.}
We propose sampling globally Lipschitz ASTs and allocating a total lag-gain
budget below one. Bounded functions such as $\tanh$, sine, and cosine allow
controlled interactions, including products of bounded factors. For example,
the proposed recurrence $0.3\tanh(x_t)\tanh(x_{t-1})+\varepsilon_{t+1}$ has
derivative bounds $0.3$ for each lag, giving a total bound of $0.6$; Gaussian
innovations satisfy the required moment condition. This route targets
stationary nonlinear dynamics on an unbounded state space. Global derivative
bounds can be conservative, particularly for polynomial feedback. A finite
direct uniform initialization can leave a finite-horizon transient; its law
must be recorded. AST sampling, certification rules, and
explanation labels are proposed adaptations in this thesis.
